% GENERATED FILE. Edit canonical lessons, then run npm run generate:books.
% canonical-source-hash: fnv1a64:55193435e500d60d
% canonical-lessons: KA-C44-look, KA-C44-listen, KA-C44-speak, KA-C44-write

\chapter{Verbs of Seeing and Saying}
\label{ch:ka-see-say}
\hlchaptermodality{44}

\begin{chapteropening}
\textbf{By the end of this chapter:} \emph{I can tell someone to look, listen, speak and write in Kannada.}

Complete KA-C44-write and say all four verbs in order.
\end{chapteropening}

\section[nōḍu]{\kn{ನೋಡು} (nōḍu) — look}

\label{lesson:KA-C44-look}



\begin{quote}
Four verbs for taking things in and putting them out, one at a time.
\end{quote}

\subsection*{You'll want to know: \kn{ನೋಡು}}
\textbf{\kn{ನೋಡು}} (\emph{nōḍu}) — \textquotedblleft{}look\textquotedblright{}.

Also 'to see' and, in the right sentence, 'to look after'. One stem covering watching, seeing and caring for is common across the Dravidian languages.

\begin{grammarlens}[title={What you've built}]
One more. Three follow, and each reuses the ones before.
\end{grammarlens}

\subsection*{Your turn}
\begin{itemize}
\raggedright
  \item \emph{Say it:} \emph{nōḍu}
  \item \emph{Say it:} it once more, slowly
  \item \emph{Recall:} say \emph{cikka}
\end{itemize}

\begin{tcolorbox}[breakable,colback=teal!4,colframe=teal!35!black,title={Before you move on}]
What does \kn{ನೋಡು} mean? (\textquotedblleft{}look\textquotedblright{}.) Say it once more.
\end{tcolorbox}
\section[kēḷu]{\kn{ಕೇಳು} (kēḷu) — listen}

\label{lesson:KA-C44-listen}



\begin{quote}
Before the new one: what did \emph{nōḍu} mean?
\end{quote}

\subsection*{You'll want to know: \kn{ಕೇಳು}}
\textbf{\kn{ಕೇಳು}} (\emph{kēḷu}) — \textquotedblleft{}listen\textquotedblright{}.

And 'to ask'. Kannada uses the same stem for taking in words and for requesting them, which is less strange than it sounds — both are aimed at what someone else knows.

\begin{grammarlens}[title={What you've built}]
One more. Three follow, and each reuses the ones before.
\end{grammarlens}

\subsection*{Your turn}
\begin{itemize}
\raggedright
  \item \emph{Say it:} \emph{kēḷu}
  \item \emph{Say it:} it once more, slowly
  \item \emph{Say it:} it after \emph{nōḍu}, so the two sit together
  \item \emph{Recall:} read \textbf{\kn{ಒಳ್ಳೆಯ}}
\end{itemize}

\begin{tcolorbox}[breakable,colback=teal!4,colframe=teal!35!black,title={Before you move on}]
What does \kn{ಕೇಳು} mean? (\textquotedblleft{}listen\textquotedblright{}.) Say it once more.
\end{tcolorbox}
\section[mātāḍu]{\kn{ಮಾತಾಡು} (mātāḍu) — speak}

\label{lesson:KA-C44-speak}



\begin{quote}
Before the new one: what did \emph{kēḷu} mean?
\end{quote}

\subsection*{You'll want to know: \kn{ಮಾತಾಡು}}
\textbf{\kn{ಮಾತಾಡು}} (\emph{mātāḍu}) — \textquotedblleft{}speak\textquotedblright{}.

Built from \emph{māta}, 'word' — itself from Sanskrit — plus \emph{āḍu}, 'to play'. To speak is to play with words.

\begin{grammarlens}[title={What you've built}]
One more. Three follow, and each reuses the ones before.
\end{grammarlens}

\subsection*{Your turn}
\begin{itemize}
\raggedright
  \item \emph{Say it:} \emph{mātāḍu}
  \item \emph{Say it:} it once more, slowly
  \item \emph{Say it:} it after \emph{kēḷu}, so the two sit together
  \item \emph{Recall:} say \emph{hosa}
\end{itemize}

\begin{tcolorbox}[breakable,colback=teal!4,colframe=teal!35!black,title={Before you move on}]
What does \kn{ಮಾತಾಡು} mean? (\textquotedblleft{}speak\textquotedblright{}.) Say it once more.
\end{tcolorbox}
\section[bare]{\kn{ಬರೆ} (bare) — write}

\label{lesson:KA-C44-write}



\begin{quote}
Before the new one: what did \emph{mātāḍu} mean?
\end{quote}

\subsection*{You'll want to know: \kn{ಬರೆ}}
\textbf{\kn{ಬರೆ}} (\emph{bare}) — \textquotedblleft{}write\textquotedblright{}.

An old Dravidian stem, unrelated to the Sanskrit words for writing. It is the same root behind \emph{baraha}, 'writing'.

\begin{grammarlens}[title={What you've built}]
That closes the run: look, listen, speak, write.
\end{grammarlens}

\subsection*{Your turn}
\begin{itemize}
\raggedright
  \item \emph{Say it:} \emph{bare}
  \item \emph{Say it:} it once more, slowly
  \item \emph{Say it:} it after \emph{mātāḍu}, so the two sit together
  \item \emph{Recall:} read \textbf{\kn{ಹಳೆಯ}}
\end{itemize}

\begin{tcolorbox}[breakable,colback=teal!4,colframe=teal!35!black,title={Before you move on}]
What does \kn{ಬರೆ} mean? (\textquotedblleft{}write\textquotedblright{}.) Say it once more.
\end{tcolorbox}
